% Generated draft: requires mathematical and source review.
\documentclass[11pt]{article}
\usepackage[margin=25mm]{geometry}
\usepackage{fontspec}
\setmainfont{Latin Modern Roman}
\usepackage{amsmath,amssymb,mathtools,unicode-math}
\setmathfont{Latin Modern Math}
\usepackage{enumitem,needspace,xurl,hyperref}
\setlength{\parindent}{0pt}
\setlength{\parskip}{5pt}
\setlength{\emergencystretch}{3em}
\newcommand{\N}{\mathbb{N}}
\newcommand{\Z}{\mathbb{Z}}
\newcommand{\Q}{\mathbb{Q}}
\newcommand{\R}{\mathbb{R}}
\newcommand{\C}{\mathbb{C}}
\newcommand{\F}{\mathbb{F}}
\newcommand{\A}{\mathbb{A}}
\newcommand{\PP}{\mathbb{P}}
\newcommand{\E}{\mathbb{E}}
\newcommand{\eps}{\varepsilon}
\newcommand{\dd}{\,\mathrm{d}}
\begin{document}
\section*{UnsolvedMath ID: 11000084}
\textbf{Source problem code:} catalogue entry
\subsection{Definitions and mathematical statement}
For a pseudo-Anosov mapping class \(f\), let \(A_f\) be its invariant Teichm\"uller geodesic axis in \(\operatorname{Teich}_g\). The quotient \(A_f/\langle f\rangle\) maps to the moduli orbifold \(\mathcal M_g\), producing the corresponding closed Teichm\"uller geodesic. A Veech curve, in the usual convention, is the image of a Teichm\"uller disk whose stabilizer is a lattice, giving a finite-area quotient of that disk.

The problem asks for an algorithm that, for each given pseudo-Anosov mapping class, decides:
\begin{enumerate}
\item whether its associated closed Teichm\"uller geodesic is simple;
\item whether that geodesic lies on a Veech curve.
\end{enumerate}
With a fixed finite generating set, a word representing a mapping class promised to be pseudo-Anosov is a natural input encoding; the source does not explicitly prescribe it. The algorithm must terminate with a decision, not merely enumerate positive instances. The source also does not specify whether simple means an embedded unparametrized image or an injective periodic parametrization. These differ for powers of a primitive mapping class, and orbifold points add a convention to be fixed.

The supplied summary confuses this question with whether the invariant foliation is a simple closed curve or multicurve. Simplicity here concerns the closed geodesic in \(\mathcal M_g\), not such a surface foliation; that reinterpretation is not adopted. The source additionally says an analogue of its Question 7.7 is unknown for hyperbolic surfaces but holds for a generic set of surfaces in \(\mathcal M_g\). Neither Question 7.7, its analogue nor genericity is defined in this record, so that dated contextual assertion cannot be expanded into an extra theorem.
\subsection{Short English statement}
Decide algorithmically whether a pseudo-Anosov's closed Teichm\"uller geodesic is simple and whether it is contained in a Veech curve, with power and orbifold conventions made explicit.
\subsection{Sources}
\begin{itemize}
\item B. Farb, chapter in \emph{Problems on Mapping Class Groups and Related Topics}, Problem 7.8, supplied PDF page 58. \url{https://www.math.uchicago.edu/~farb/papers/mcgbook.pdf}. Recorded access 2026-07-29; extraction pdf-font-structure. The cited Question 7.7 is absent; no separate algorithm or Veech reference is supplied.
\item Local UnsolvedMath registry, \texttt{data/unsolvedmath/problems.json}, permanent ID \(11000084\), exact source code \texttt{catalogue entry}; supplied OPEN-TRIAGE summary reports unsettled decidability but gives an incompatible foliation description and no primary citation.
\end{itemize}
\end{document}
